% BEGIN VOL08-EXPANSION VIII32
\section{Supplementary worked examples}
\label{sec:viii32-pedagogy}

\begin{example}[Euler class of the tangent bundle of S2]\label{ex:viii32-ped-01}
For the oriented tangent bundle $TS^2$,
\[
\langle e(TS^2),[S^2]\rangle=\chi(S^2)=2.
\]
Therefore $e(TS^2)\ne0$, so $TS^2$ admits no nowhere-zero section: the hairy-ball theorem follows immediately.
\end{example}

\begin{example}[Trivial bundle and a nowhere-zero section]\label{ex:viii32-ped-02}
If $E=B\times\mathbb R^r$ has positive rank, the constant section $s(x)=(x,e_1)$ never vanishes. Consequently $e(E)=0$. The vanishing is necessary for a nowhere-zero section, though in general it need not be sufficient without dimension or rank hypotheses.
\end{example}

\begin{example}[Sphere bundle as the boundary of the disk bundle]\label{ex:viii32-ped-03}
Choosing a bundle metric gives the disk bundle $D(E)$ and sphere bundle $S(E)=\partial D(E)$. The pair $(D(E),S(E))$ carries the Thom class, while the sphere bundle records directions of nonzero vectors; a nowhere-zero section of $E$ is equivalent after normalization to a section of $S(E)$.
\end{example}

\section{Graded supplementary exercises}
The following exercises progress from direct computation and construction to proof, hypothesis testing, application, and synthesis.

\subsection*{Standard computations and constructions}

\begin{exercise}[Sphere fiber]\label{exr:viii32-ped-01}
What is the fiber of the sphere bundle of a rank-$r$ real vector bundle?
\end{exercise}
\begin{hint}
Normalize nonzero vectors.
\end{hint}
\begin{solution}
The fiber is $S^{r-1}$.
\end{solution}

\begin{exercise}[Normalize section]\label{exr:viii32-ped-02}
How does a nowhere-zero section of $E$ produce a section of $S(E)$ after choosing a metric?
\end{exercise}
\begin{hint}
Divide by its norm.
\end{hint}
\begin{solution}
$x\mapsto s(x)/\|s(x)\|$ is a sphere-bundle section.
\end{solution}

\begin{exercise}[Euler degree]\label{exr:viii32-ped-03}
In what degree does $e(E)$ lie for an oriented rank-$r$ real bundle?
\end{exercise}
\begin{hint}
It is the zero-section pullback of the degree-$r$ Thom class.
\end{hint}
\begin{solution}
$e(E)\in H^r(B;\mathbb Z)$.
\end{solution}

\begin{exercise}[Trivial Euler class]\label{exr:viii32-ped-04}
Compute $e(B\times\mathbb R^r)$ for $r>0$.
\end{exercise}
\begin{hint}
Use a constant nonzero section.
\end{hint}
\begin{solution}
It is zero.
\end{solution}

\begin{exercise}[TS2 Euler number]\label{exr:viii32-ped-05}
Compute $\langle e(TS^2),[S^2]\rangle$.
\end{exercise}
\begin{hint}
Use the Euler characteristic.
\end{hint}
\begin{solution}
It equals $2$.
\end{solution}

\subsection*{Proofs}

\begin{exercise}[Section forces vanishing]\label{exr:viii32-ped-06}
Prove a nowhere-zero section forces $e(E)=0$ for an oriented bundle.
\end{exercise}
\begin{hint}
Move the zero section along the nonzero section or use disjointness.
\end{hint}
\begin{solution}
The section is homotopic through bundle sections to a section disjoint from the zero section, so its intersection with the zero section vanishes; equivalently the pullback of the Thom class along a nonzero section is zero, hence $e(E)=0$.
\end{solution}

\begin{exercise}[Euler naturality]\label{exr:viii32-ped-07}
Show $e(f^*E)=f^*e(E)$ for orientation-preserving pullback.
\end{exercise}
\begin{hint}
Use naturality of the Thom class and zero section.
\end{hint}
\begin{solution}
Pullback commutes with the Thom class and with the zero-section square, giving the stated equality.
\end{solution}

\begin{exercise}[Euler number and zeros]\label{exr:viii32-ped-08}
Explain why a transverse section of an oriented rank-$n$ bundle over a closed oriented $n$-manifold has signed zero count equal to $\langle e(E),[M]\rangle$.
\end{exercise}
\begin{hint}
Each zero is an oriented local intersection with the zero section.
\end{hint}
\begin{solution}
The pullback Thom class localizes at transverse zeros; summing local intersection signs evaluates the Euler class on the fundamental class.
\end{solution}

\begin{exercise}[Hairy ball]\label{exr:viii32-ped-09}
Deduce that $S^2$ has no nowhere-zero tangent vector field.
\end{exercise}
\begin{hint}
Use the nonzero Euler number.
\end{hint}
\begin{solution}
A nowhere-zero tangent field would force $e(TS^2)=0$, contradicting evaluation $2$.
\end{solution}

\subsection*{Counterexamples and hypothesis tests}

\begin{exercise}[Euler zero implies section always]\label{exr:viii32-ped-10}
Test: $e(E)=0$ always guarantees a nowhere-zero section.
\end{exercise}
\begin{hint}
Euler class is only the primary obstruction in general.
\end{hint}
\begin{solution}
False without further hypotheses; higher obstructions may remain.
\end{solution}

\begin{exercise}[Sphere bundle always trivial]\label{exr:viii32-ped-11}
Test: the unit sphere bundle $S(E)$ is a product $B\times S^{r-1}$.
\end{exercise}
\begin{hint}
Nontrivial vector bundles have nontrivial sphere bundles.
\end{hint}
\begin{solution}
False. For example, sphere bundles associated to nontrivial tangent bundles need not be products.
\end{solution}

\begin{exercise}[Odd sphere tangent Euler]\label{exr:viii32-ped-12}
Test: $\langle e(TS^{2k+1}),[S^{2k+1}]\rangle=2$.
\end{exercise}
\begin{hint}
Use $\chi(S^{2k+1})=0$.
\end{hint}
\begin{solution}
False; the Euler number is zero for odd-dimensional spheres.
\end{solution}

\subsection*{Applications and investigations}

\begin{exercise}[Vector-field diagnostics]\label{exr:viii32-ped-13}
How can computing an Euler class guide a search for global nonvanishing fields?
\end{exercise}
\begin{hint}
Nonzero Euler class is a certificate of impossibility.
\end{hint}
\begin{solution}
If the Euler class is nonzero, no algorithm can produce a global nowhere-zero section without violating continuity; vanishing removes this primary obstruction.
\end{solution}

\begin{exercise}[Sphere-bundle obstruction]\label{exr:viii32-ped-14}
Why can the existence of a section of $S(E)$ be formulated as an obstruction problem?
\end{exercise}
\begin{hint}
Build the section over skeleta.
\end{hint}
\begin{solution}
One extends a partial sphere-bundle section skeleton by skeleton; obstruction classes measure failures to extend, with the Euler class appearing as the primary obstruction in the relevant rank.
\end{solution}

\subsection*{Challenge problems}

\begin{exercise}[Tangent surface]\label{exr:viii32-ped-15}
For a closed oriented surface $\Sigma_g$, compute the Euler number of $T\Sigma_g$.
\end{exercise}
\begin{hint}
Use Gauss--Bonnet/Poincare--Hopf.
\end{hint}
\begin{solution}
$\langle e(T\Sigma_g),[\Sigma_g]\rangle=\chi(\Sigma_g)=2-2g$.
\end{solution}

\begin{exercise}[Euler synthesis]\label{exr:viii32-ped-16}
Connect sphere-bundle sections, Thom classes, zeros of sections, and Euler characteristic.
\end{exercise}
\begin{hint}
Normalize nonzero sections and compare with the zero section.
\end{hint}
\begin{solution}
A nowhere-zero section is a sphere-bundle section; the Thom class detects intersection with the zero section; its pullback is the Euler class; for tangent bundles its evaluation equals the Euler characteristic, converting topology into a zero-count obstruction.
\end{solution}

% END VOL08-EXPANSION VIII32
